% ============================================================
%  cap01.tex  --  Capítulo 1: Regimes de Capitalização
%  (Conteúdo reconstituído da Aula 1 -- não enviado como PDF;
%   base: Bueno, Rangel & Santos, Cap. 1)
% ============================================================

\chapter{Regimes de Capitalização}

\section{Conceito de Juros}

\begin{definicao}
\textbf{Juros} são a remuneração pelo sacrifício do consumo presente, ou equivalentemente, o custo de transferir poder de compra ao longo do tempo.
Se um agente empresta o capital $P$ (principal) por $n$ períodos a uma taxa $i$ por período, o montante recebido ao final é $S = P + J$, onde $J$ representa os juros totais.
\end{definicao}

A forma como os juros se acumulam ao longo do tempo define o \textbf{regime de capitalização}. Os três regimes principais são:
\begin{enumerate}[label=\textbf{\arabic*.}]
  \item \textbf{Juros simples} -- os juros são calculados sempre sobre o principal $P$;
  \item \textbf{Juros compostos} -- os juros de cada período se incorporam ao principal (``juros sobre juros'');
  \item \textbf{Capitalização contínua (instantânea)} -- limite do regime composto quando o intervalo de capitalização tende a zero.
\end{enumerate}

% ─────────────────────────────────────────────────────────────────────────────
\section{Juros Simples}

No regime de \textbf{juros simples}, a cada período acrescentam-se juros iguais a $i \cdot P$:
\[
  J_t = i \cdot P, \qquad t = 1, 2, \ldots, n.
\]
O montante após $n$ períodos é portanto:
\begin{resultado}
\[
  S = P\,(1 + i\,n).
\]
\end{resultado}

\subsection*{Convenções de contagem de dias}

Na prática brasileira, dois padrões coexistem:

\begin{center}
\renewcommand{\arraystretch}{1.3}
\begin{tabular}{lll}
\toprule
\textbf{Convenção} & \textbf{Ano} & \textbf{Mês} \\
\midrule
Dias corridos & 360 dias & 30 dias \\
Dias úteis    & 252 dias úteis & 21 dias úteis \\
\bottomrule
\end{tabular}
\end{center}

\begin{nota}
Devido a feriados e ao calendário, nem todos os meses efetivos têm o mesmo número de dias úteis. A estratégia correta é sempre usar os dias \emph{exatos} do período analisado.

\smallskip
Ferramentas práticas:
\begin{itemize}[leftmargin=1.5em]
  \item \textbf{HP 12C}: calcula dias efetivos entre duas datas no calendário de 360 ou 365 dias (\texttt{g} $\Delta$DYS).
  \item \textbf{Excel}: função \texttt{DIAS360} (corridos) ou \texttt{DIATRABALHOTOTAL} (úteis, requer tabela de feriados).
\end{itemize}
\end{nota}

\begin{exemplo}{Comparação de rendimentos em juros simples}
Um título paga juros simples de $2{,}5\%$ ao mês. Qual o montante de R\$\,10.000 investidos por 45 dias corridos?

\medskip
$n = 45/30 = 1{,}5$ meses, portanto:
\[
  S = 10{.}000\,(1 + 0{,}025 \times 1{,}5) = 10{.}000 \times 1{,}0375 = \text{R\$\,}10.375{,}00.
\]
\end{exemplo}

% ─────────────────────────────────────────────────────────────────────────────
\section{Juros Compostos}

No regime de \textbf{juros compostos}, ao fim de cada período os juros são incorporados ao saldo, tornando-se base de cálculo do período seguinte:
\[
  S_1 = P(1+i),\quad S_2 = S_1(1+i) = P(1+i)^2,\quad \ldots
\]

\begin{resultado}
\[
  \boxed{S = P\,(1+i)^n}
\]
\end{resultado}

\subsection{Equivalência de taxas}

Duas taxas $i_1$ (para o período $n_1$) e $i_2$ (para o período $n_2$) são \textbf{equivalentes} sob juros compostos se produzem o mesmo montante:
\[
  (1+i_1)^{n_1} = (1+i_2)^{n_2}
  \quad\Rightarrow\quad
  i_2 = (1+i_1)^{n_1/n_2} - 1.
\]

\begin{exemplo}{Equivalência mensal--anual}
A taxa de $12\%$ a.a. equivale a que taxa mensal?
\[
  i_m = (1{,}12)^{1/12} - 1 \approx 0{,}9489\% \text{ a.m.}
\]
(Note que a taxa proporcional de $12\%/12 = 1\%$ a.m. é \emph{diferente}; a proporcionalidade é válida apenas em juros simples.)
\end{exemplo}

\subsection{Taxa no período}

Para um investimento de $n$ períodos à taxa $i$ por período, a \textbf{taxa total no período} (ou taxa efetiva do período) é:
\[
  i_{\text{per}} = (1+i)^n - 1.
\]

\begin{exemplo}{Rendimento acumulado em 3 anos}
Um fundo rende $0{,}9\%$ a.m. Qual o rendimento acumulado em 36 meses?
\[
  i_{36} = (1{,}009)^{36} - 1 \approx 38{,}02\%.
\]
\end{exemplo}

\subsection{Taxa de juros variável}

Quando a taxa varia a cada período $t$, o montante é:
\begin{resultado}
\[
  S = P \prod_{t=1}^{n}(1+i_t).
\]
\end{resultado}

\begin{exemplo}{Taxa variável -- 3 semestres}
Taxas semestrais: $5\%$, $4{,}5\%$, $6\%$. Montante de R\$\,1.000:
\[
  S = 1{.}000 \times 1{,}05 \times 1{,}045 \times 1{,}06 \approx \text{R\$\,}1.163{,}78.
\]
\end{exemplo}

% ─────────────────────────────────────────────────────────────────────────────
\section{Capitalização Contínua (Instantânea)}

À medida que o intervalo de capitalização diminui para zero, o montante converge para:
\begin{resultado}
\[
  \boxed{S = P\,e^{r\,n}}
\]
\end{resultado}
onde $r$ é a \textbf{taxa contínua} (ou taxa instantânea de juros).

\subsection*{Relação entre taxa contínua e taxa discreta}

Igualando os montantes dos dois regimes:
\[
  P\,e^{r\,n} = P\,(1+i)^n
  \quad\Rightarrow\quad
  e^r = 1+i
  \quad\Rightarrow\quad
  \boxed{r = \ln(1+i)}
  \quad\Leftrightarrow\quad
  \boxed{i = e^r - 1.}
\]

\begin{exemplo}{Conversão taxa discreta $\leftrightarrow$ contínua}
\begin{itemize}
  \item $i = 12\%$ a.a. $\Rightarrow$ $r = \ln(1{,}12) \approx 11{,}33\%$ a.a. (contínua).
  \item $r = 10\%$ a.a. (contínua) $\Rightarrow$ $i = e^{0{,}10}-1 \approx 10{,}52\%$ a.a. (discreta).
\end{itemize}
\end{exemplo}

\begin{moderno}[title={Taxa Selic em capitalização contínua (2024)}]
A taxa Selic fixada em $10{,}50\%$ a.a. (discreta, base 252 d.u.) corresponde à taxa contínua:
\[
  r = \ln(1{,}1050) \approx 9{,}98\% \text{ a.a.}
\]
Modelos de apreçamento de derivativos (Black--Scholes, HJM) tipicamente trabalham com $r$ contínuo.
\end{moderno}

% ─────────────────────────────────────────────────────────────────────────────
\section{Comparação dos Regimes}

Para um mesmo capital $P$, taxa $i$ (ou taxa contínua equivalente $r = \ln(1+i)$) e prazo $n$:
\begin{center}
\renewcommand{\arraystretch}{1.4}
\begin{tabular}{llll}
\toprule
\textbf{Regime}     & \textbf{Montante}    & \textbf{Cresce} & \textbf{Uso típico} \\
\midrule
Juros simples       & $P(1+in)$            & Linearmente     & Títulos curto prazo, CDI/dia \\
Juros compostos     & $P(1+i)^n$           & Exponencialmente & Maioria das operações \\
Capitalização cont. & $Pe^{rn}$            & Exponencialmente & Finanças quantitativas \\
\bottomrule
\end{tabular}
\end{center}

\begin{nota}
Para $n < 1$ e mesma taxa: juros simples $>$ compostos. Para $n > 1$: compostos $>$ simples. Na capitalização contínua e composta, os montantes coincidem quando $r = \ln(1+i)$.
\end{nota}
